\section{Empirical Results}
\label{sec:results}

\subsection{Monolingual Retrieval Baselines}
\label{subsec:mono_results}

Top-1 retrieval accuracy on the C++ and Java monolingual cohorts under the Phase~1 weight configuration ($w_{BL}=0.35, w_{WL}=0.40, w_{CES}=0.25$) with the asymmetric Tversky index is reported in Table~\ref{tab:mono_results}. Without parameter retuning, the entire pipeline—which includes context normalization and STL injection—achieves 90.00\% on C++ ($n=360/400$) and 85.75\% on Java ($n=343/400$). The 23.25\% zero-CES failure rate in the uncurated ablation set (Table~\ref{tab:ces_failure}) and Java's intrinsically higher CES extraction noise from object-oriented verbosity are the two compounding causes that contribute to the 4.25 pp difference between C++ and Java monolingual accuracy. $91.00\% \pm 3.3\%$ is the result of five-fold cross-validation on the C++ cohort; $92.00\% \pm 2.92\%$ is the result of Java five-fold cross-validation on the structurally compatible {subset. Because} the global figure includes the uncurated zero-CES failure penalty while cross-validation is only calculated on compliant programs, the Java CV surpasses the 85.75\% global baseline. On a workstation running Joern v1.1.1 in standalone mode with an Intel Core i5-10400 CPU (2.90 GHz, 6 cores) and 16 GB of RAM, processing time averaged ${\sim}1.2$ seconds per submission. 

\begin{table}[ht] 
\caption{Monolingual Retrieval Accuracy. Phase~1 weights ($w_{BL}=0.35, w_{WL}=0.40, w_{CES}=0.25)$ with an asymmetric Tversky index ($\alpha=0.1$, $\beta=0.9$) are used in the results. The final, perfectly balanced 400-program assessment dataset is reflected in the Java monolingual accuracy (85.75\%); the discordant cells of Table~\ref{tab:mcnemar_contingency} are used to calculate the McNemar odds ratio (OR$=b/c=4.875$).}
\label{tab:mono_results}
\centering
\resizebox{\columnwidth}{!}{%
\begin{tabular}{|l|c|c|c|}
\hline
\textbf{Dataset} & \textbf{Correct/Total} & \textbf{Top-1 Acc.} & \textbf{95\% Wilson CI} \\
\hline
C++ Monolingual                   & 360/400 & 90.00\% & [86.73\%, 92.57\%] \\
Java Monolingual                  & 343/400 & 85.75\% & [82.10\%, 88.79\%] \\
\hline
C++ 5-fold CV                     & ---     & 91.00\% & [$\pm$3.3\%] \\
Java 5-fold CV (compliant subset) & ---     & 92.00\% & [$\pm$2.92\%] \\
\hline
\end{tabular}%
}
\end{table}

\subsection{Cross-Language Retrieval}
\label{subsec:cross_lang_results}

The framework was evaluated on the 400-program cross-language dataset spanning C, C++, and Java under the Context Normalization Protocol, asymmetric Tversky similarity, and the strict hard-filtering constraint (Section~\ref{subsec:late_fusion}). Input submissions are balanced across languages: 130 C, 130 C++, and 140 Java programs across 20 problems.

\textbf{Directional composition disclosure:} The cross-language evaluation exhibits a pronounced directional imbalance in top-1 retrieval \textit{outcomes}, arising from the shared \texttt{c2cpg} parser frontend used for both C and C++. Because C and C++ programs produce nearly identical CPG node-type vocabularies, their WL hashes are structurally very close; C/C++ queries therefore almost always rank another C/C++ reference above any Java reference in the competitive top-1 race. Consequently, 246 of 400 cross-language query-reference outcome pairs (61.5\%) are C$\leftrightarrow$C++ (shared-frontend) pairs, while Java$\rightarrow$C/C++ accounts for $n=140$ (35.0\%). The C$\rightarrow$Java ($n=11$) and C++$\rightarrow$Java ($n=3$) cells are statistically underpowered for directional inference; while included in the overall full-corpus aggregate of 87.25\% ($349/400$), they are also reported separately alongside the core-direction aggregate ($342/386 = 88.60\%$) in Table~\ref{tab:directional_results} for complete transparency. 

\begin{table}[ht]
\caption{Cross-Language Top-1 Retrieval Accuracy by Directional Pair. Inputs for all languages are balanced (130~C, 130~C++, 140~Java). Outcomes per direction indicate competitive interference due to the use of a common C/C++ parser frontend. $^{\dagger}$The C$\rightarrow$Java ($n=11$) and C++$\rightarrow$Java ($n=3$) directions are statistically underpowered but retained for complete cross-language reporting; the aggregate accuracy of 87.25\% ($349/400$) represents the full-corpus metric across all directional pairs, while the core-direction aggregate (Java$\rightarrow$C/C++ and C$\leftrightarrow$C++) yields 88.60\% ($342/386$).}
\label{tab:directional_results}
\centering
\resizebox{\columnwidth}{!}{%
\begin{tabular}{|l|c|c|c|}
\hline
\textbf{Direction} & \textbf{Correct/Total} & \textbf{Top-1 Acc.} & \textbf{95\% Wilson CI} \\
\hline
Java$\rightarrow$C/C++        & 128/140  & 91.43\%          & [85.60\%, 95.00\%] \\
C$\leftrightarrow$C++         & 214/246  & 86.99\%          & [82.21\%, 90.63\%] \\
C$\rightarrow$Java$^{\dagger}$    &   6/11   & 54.55\%          & [26.67\%, 80.12\%] \\
C++$\rightarrow$Java$^{\dagger}$  &   1/3    & 33.33\%          & [3.34\%, 77.72\%]  \\
\hline
\textbf{Core-Direction Subtotal} & \textbf{342/386} & \textbf{88.60\%} & \textbf{[85.04\%, 91.46\%]} \\
\textbf{Full-Corpus Aggregate}   & \textbf{349/400} & \textbf{87.25\%} & \textbf{[83.60\%, 90.20\%]} \\
\hline
\end{tabular}%
}
\end{table}

\textbf{Aggregate interpretation:} The 87.25\% full-corpus cross-language accuracy ($349/400$) captures all 400 evaluated submissions across the six directional pairs. When restricted to the primary, statistically robust directions (Java$\rightarrow$C/C++ and C$\leftrightarrow$C++, $n=386$), the core-direction aggregate reaches 88.60\% ($342/386$). The sparse C$\rightarrow$Java and C++$\rightarrow$Java directions ($n=11$ and $n=3$) reflect frontend competitive displacement rather than strategy misclassification. This structural asymmetry is a critical caveat; the aggregate performance should be interpreted as frontend-stratified and not as a claim of symmetric, bidirectional cross-language generalization.

\textbf{Analysis of directional accuracy and similarity scores:} A straightforward examination of the resultant combined similarity matrix lends empirical evidence to support the hypothesis that the 61.5\% C$\leftrightarrow$C++ share in retrieval results indicates competitive displacement by the shared \texttt{c2cpg} parser frontend rather than the failure of Java program retrieval. As indicated in Table~\ref{tab:displacement}, accuracy in each of the directions and the mean top-1 similarity scores have been computed for the 400 evaluations. While the Java$\rightarrow$C/C++ direction attains a remarkably high \textbf{top-1 accuracy of 91.43\%} [95\% Wilson Confidence Interval (CI): 85.60\%, 95.00\%] on $n=140$ pairs---\textit{higher than the C$\leftrightarrow$C++ direction} (86.99\% [95\% Wilson CI: 82.21\%, 90.63\%], $n=246$ pairs)---it generates substantially low similarity scores compared to C$\leftrightarrow$C++ (mean: 0.796 vs.\ 0.911). Such a discrepancy negates the possibility of interpreting the sparsity of Java retrieval results in terms of Java language retrieval failure. Moreover, the average difference between the winning score and the correct-reference score in Java$\rightarrow$C/C++ translations is $\Delta=0.0037$ (median: 0.0000), while 96.4\% of Java$\rightarrow$C/C++ retrievals have a gap smaller than 0.02---suggesting that the retrieval threshold in this direction is rather tight and the model is strongly discriminative here. The sparsity in C$\rightarrow$Java and C++$\rightarrow$Java translations can be attributed to competitive displacement: the fact that structurally identical C/C++ reference is higher-ranked than the correct Java reference at top-1 position is due to the fact that the two references use the same node vocabulary \texttt{c2cpg}, but does not mean that the latter is misranked semantically. This observation is supported by Table~\ref{tab:displacement}. It should be noted that due to the nature of our experiment setup, competitive displacement cannot be established as causal relationship since it is based on the observations only; nevertheless, we consider this explanation very plausible. Each retrieval instance was done from about 6--7 references per query per language; implications are discussed in Section~\ref{subsec:weight_stability_discussion}.
\begin{table*}[!t]
\caption{Directional Accuracy and Similarity Score Analysis. Computed from the final combined similarity matrix (Phase~1 weights, Tversky index, strict hard-filtering). ``Mean winner score'' is the mean top-1 similarity across all queries in that direction; ``Mean correct-ref score'' is the mean similarity to the ground-truth correct reference. ``Mean gap'' is winner $-$ correct-ref score. The Java$\rightarrow$C/C++ direction achieves \textit{higher} retrieval accuracy than C$\leftrightarrow$C++ despite lower raw similarity scores, providing observational support for the competitive displacement hypothesis.}
\label{tab:displacement}
\centering
\begin{tabular}{|l|c|c|c|c|c|}
\hline
\textbf{Direction} & \textbf{n (\% of 400)} & \textbf{Top-1 Acc.} & \textbf{Mean Winner Score} & \textbf{Mean Correct-Ref Score} & \textbf{Mean Gap ($\Delta$)} \\
\hline
Java$\rightarrow$C/C++ & 140 (35.0\%) & \textbf{91.43\%} & 0.796 & 0.792 & 0.004 \\
C$\leftrightarrow$C++  & 246 (61.5\%) & 86.99\%          & 0.911 & 0.900 & 0.010 \\
\hline
\end{tabular}
\end{table*}

\textbf{Single-view ablation study:} To isolate the individual retrieval contribution of each representation perspective, Table~\ref{tab:single_view_ablation} reports standalone top-1 retrieval performance on the cross-language dataset ($N=400$) when setting two of the three late-fusion weights to zero. While individual views achieve between 71.68\% and 79.50\% retrieval accuracy, late fusion of all three views $(w_{\text{BL}}=0.35, w_{\text{WL}}=0.40, w_{\text{CES}}=0.25)$ achieves 87.25\%, demonstrating a +7.75~pp gain over the strongest standalone view (lexical BL).

\begin{table}[ht]
\caption{Single-View Ablation Study on Cross-Language Dataset ($N=400$, 95\% Wilson CIs under strict hard-filtering). Each isolated view sets the other two view weights to zero.}
\label{tab:single_view_ablation}
\centering
\begin{tabular}{|l|c|c|c|c|}
\hline
\textbf{Representation View} & $\boldsymbol{w_{\text{BL}}}$ & $\boldsymbol{w_{\text{WL}}}$ & $\boldsymbol{w_{\text{CES}}}$ & \textbf{Top-1 Acc. [95\% CI]} \\
\hline
BL-only (Lexical view)   & 1.00 & 0.00 & 0.00 & 79.50\% [75.35\%, 83.11\%] \\
WL-only (Structural view)& 0.00 & 1.00 & 0.00 & 71.68\% [67.07\%, 75.88\%] \\
CES-only (Semantic view) & 0.00 & 0.00 & 1.00 & 72.68\% [68.11\%, 76.82\%] \\
\hline
\textbf{Fused Framework} & \textbf{0.35} & \textbf{0.40} & \textbf{0.25} & \textbf{87.25\% [83.60\%, 90.20\%]} \\
\hline
\end{tabular}
\end{table}

\textbf{Baseline comparison:} Table~\ref{tab:baseline_comparison} shows our proposed framework versus external baselines on an identical 400-program cross-language data set. Non-semantic code-clone detectors (JPlag, MOSS) and pretrained neural encoders (CodeBERT, GraphCodeBERT, UniXcoder) are presented alongside the internal BL-view ablation ($w_{\text{BL}}=1.0$) to distinguish external comparative baselines from internal component ablations.

\begin{table}[ht]
\caption{Baseline Comparison on Cross-Language Dataset ($N=400$, under identical hard-filtering). Neural models evaluated zero-shot (CLS-token cosine similarity, no task-specific fine-tuning) unless otherwise noted. GraphCodeBERT evaluated on raw code tokens \textit{without} DFG preprocessing; this result reflects degraded zero-shot performance and should be interpreted as a lower bound on GraphCodeBERT capability (see Section~\ref{subsec:neural_comparison}).}
\label{tab:baseline_comparison}
\centering
\begin{tabular}{|l|c|c|}
\hline
\textbf{System} & \textbf{Top-1 Accuracy} & \textbf{95\% Wilson CI} \\
\hline
\multicolumn{3}{|c|}{\cellcolor{gray!20}\textbf{Non-Semantic Lower Bounds}} \\
\hline
JPlag (token-fingerprinting)  & 51.50\% & [46.73\%, 56.26\%] \\
MOSS (token-fingerprinting)   & 54.20\% & [49.43\%, 58.93\%] \\
\hline
\multicolumn{3}{|c|}{\cellcolor{gray!20}\textbf{Neural & External IR Baselines}} \\
\hline
GraphCodeBERT (zero-shot)$^{\dagger}$         & 58.90\% & [54.01\%, 63.62\%] \\
CodeBERT (zero-shot)$^{\dagger}$               & 59.50\% & [54.72\%, 64.10\%] \\
CodeBERT (fine-tuned, $n{=}80$)$^{\S}$        & 66.25\% & [55.36\%, 75.65\%] \\
UniXcoder (zero-shot)$^{\dagger}$             & 68.92\% & [64.22\%, 73.26\%] \\
UniXcoder (fine-tuned, $n{=}80$)$^{\ddagger}$ & 71.37\% & [62.15\%, 79.12\%] \\
\hline
\multicolumn{3}{|c|}{\cellcolor{gray!20}\textbf{Framework Ablation \& Fused Model}} \\
\hline
Framework BL-View Ablation ($w_{\text{BL}}{=}1$)  & 79.50\% & [75.35\%, 83.11\%] \\
\textbf{Proposed Fused Framework}               & \textbf{87.25\%} & \textbf{[83.60\%, 90.20\%]} \\
\hline
\end{tabular}
\begin{flushleft}
\footnotesize
$^{\dagger}$Zero-shot: Cosine similarity based on CLS-token, $N=400$, no fine-tuning (GraphCodeBERT zero-shot positioning is discussed in Section~\ref{subsec:neural_comparison}).\\
$^{\ddagger}$UniXcoder fine-tuning protocol is detailed in Section~\ref{subsec:neural_finetuning}. Framework also fine-tuned on same 80 programs: 88.75\% [80.49\%, 93.83\%].\\
$^{\S}$CodeBERT fine-tuned on same held-out 80 programs; InfoNCE objective function, AdamW with learning rate $lr{=}2{\times}10^{-5}$ for 3 epochs.
\end{flushleft}
\end{table}

\textbf{Statistical analysis:} The proposed fused framework yielded 31 additional correct retrievals over the standalone BL-view ablation (349 vs.\ 318 correct out of 400). Table~\ref{tab:mcnemar_contingency} provides the full 2$\times$2 contingency matrix; cell counts $b=39$ and $c=8$ are the discordant pairs driving the McNemar statistic. This paired test quantifies the incremental statistical benefit of fusing structural (WL) and semantic (CES) views over lexical TF-IDF alone.

\begin{table}[ht]
\caption{McNemar Contingency Matrix: Proposed Fused Framework vs.\ Standalone BL-View Ablation ($N=400$ cross-language submissions). Discordant cells: $b=39$ (framework correct, BL wrong), $c=8$ (framework wrong, BL correct).}
\label{tab:mcnemar_contingency}
\centering
\begin{tabular}{|l|c|c|}
\hline
 & \textbf{BL Ablation Correct} & \textbf{BL Ablation Wrong} \\
\hline
\textbf{Fused Framework Correct} & $a = 310$ & $b = 39$ \\
\textbf{Fused Framework Wrong}   & $c = 8$   & $d = 43$ \\
\hline
\end{tabular}
\end{table}

From cells $b=39$ and $c=8$: raw McNemar~\cite{ref_mcnemar1947} $\chi^2 = (b-c)^2/(b+c) = (31)^2/47 \approx 20.45$ ($p < 0.0001$); Yates-corrected: $\chi^2 = (|b-c|-1)^2/(b+c) = (30)^2/47 \approx 19.15$ ($p < 0.0001$, one degree of freedom). For a paired McNemar design, the appropriate odds ratio is the ratio of discordant cells: OR $= b/c = 39/8 = 4.875$ [95\% CI: 2.26, 10.52], computed as $\exp(\ln(b/c) \pm 1.96\sqrt{1/b+1/c})$~\cite{ref_fleiss2003}. Beyond statistical significance, the 7.75 percentage point improvement demonstrates the substantial gain achieved by combining structural graph kernels and semantic signatures with lexical features.

\begin{figure}[ht]
    \centering
    \includegraphics[width=\linewidth]{fig2.png}
    \caption{CES weight sensitivity across evaluation contexts. $w_{CES}$ varies along the x-axis while $(w_{BL}, w_{WL})$ maintain their Phase~1 ratios. C++ exhibits a large plateau range, which is expected, given the high coverage of CES extraction. Java shows a large decline in accuracy when CES weight becomes large, which is again due to OO verbosity making CES extraction harder; the cross-language method behaves like the Java one because the same reason holds for cross-language submissions. All methods perform above 85\% accuracy in the Phase~1 configuration, where $w_{CES}=0.25$. Each data point comes from a full population deterministic test, hence no model variance is present. Uncertainty stems from annotation accuracy ($\kappa=0.86$) and problem space sampling, not stochastic model variance.}
    \label{fig:weight_sensitivity}
\end{figure}

\subsection{Per-Problem Accuracy Breakdown}
\label{subsec:per_problem}

Table~\ref{tab:per_problem} shows the top-1 retrieval accuracy for each problem type on the cross-language dataset. The aggregate 87.25\% conceals meaningful variation attributable to strategy diversity and taxonomy coverage gaps rather than language-level failure.

\begin{table*}[!t]
\caption{Per-Problem Top-1 Retrieval Accuracy on the Cross-Language Dataset ($N=20$ submissions per problem, balanced across C/C++/Java). Problems with accuracy below 80\% are highlighted. ``Primary Failure Lang.'' indicates the language context contributing the majority of retrieval errors; ``Dominant Src Lang.'' indicates the dominant query source language among the cross-language pairs for that problem.}
\label{tab:per_problem}
\centering
\begin{tabular}{|c|l|c|c|l|l|}
\hline
\textbf{ID} & \textbf{Problem} & \textbf{Correct/Total} & \textbf{Acc.} & \textbf{Primary Failure Lang.} & \textbf{Dominant Src Lang.} \\
\hline
P1  & Sum of Array                & 20/20 & 100\% & ---                          & Mixed \\
P2  & Max/Min Element             & 20/20 & 100\% & ---                          & Mixed \\
P3  & Linear Search               & 19/20 & 95\%  & Mixed                        & Mixed \\
P4  & Bubble Sort                 & 19/20 & 95\%  & Mixed                        & Mixed \\
P5  & Palindrome Check            & 19/20 & 95\%  & Mixed                        & Mixed \\
P6  & String Reversal             & 18/20 & 90\%  & Mixed                        & Mixed \\
P7  & Fibonacci (Iterative)       & 18/20 & 90\%  & Mixed                        & Mixed \\
P8  & GCD (Euclidean, recursive)  & 17/20 & 85\%  & Mixed                        & Mixed \\
P9  & Factorial (recursive)       & 18/20 & 90\%  & Mixed                        & Mixed \\
P10 & Prime Check                 & 18/20 & 90\%  & Mixed                        & Mixed \\
P11 & Binary Search               & 17/20 & 85\%  & Mixed                        & Mixed \\
P12 & Insertion Sort              & 17/20 & 85\%  & Mixed                        & Mixed \\
P13 & Selection Sort              & 18/20 & 90\%  & Mixed                        & Mixed \\
% \rowcolor{yellow!30}
P14 & Array Rotation              & 15/20 & 75\%  & C/C++ (reversal variant)     & C/C++ dominant \\
% \rowcolor{yellow!30}
P15 & Common Element Detection    & 14/20 & 70\%  & Java (hash-set idiom)        & Java dominant \\
P16 & Digit Sum                   & 17/20 & 85\%  & Mixed                        & Mixed \\
P17 & Power Computation           & 17/20 & 85\%  & Mixed                        & Mixed \\
P18 & Matrix Addition             & 16/20 & 80\%  & Mixed                        & Mixed \\
P19 & LCM                         & 16/20 & 80\%  & Mixed                        & Mixed \\
P20 & Fibonacci (Recursive)       & 17/20 & 85\%  & Mixed                        & Mixed \\
\hline
\end{tabular}
\end{table*}

\textbf{P14 -- Array Rotation (75\%):} Five instances of retrieval failures appear to be due to the limitations on strategy diversity and taxonomy coverage. In the case of array rotation, the reversal-based strategy (a 3-reversal operation which creates the following patterns: CONDITIONAL\_SWAP and NARROWING\_WINDOW) is used together with an index shift approach. Though both strategies are correct, their vectors differ greatly; therefore, while the reference database mostly contains solutions based on indexes, reversal-based submissions are mislabeled. This pattern is consistent with a taxonomy-coverage artifact rather than a cross-language failure.

\textbf{P15 -- Common Element Detection (70\%):} Six failed retrievals demonstrate strategy diversity and taxonomy coverage limitations among three implementations: nested loops (O($n^2$) brute force approach), merge-scan using a prior sort (O($n \log n$)), and Java's hash-set check (\texttt{HashSet.contains()}). Each of these approaches leads to a significantly different CES vector, and the 20-program reference set is distributed unevenly among these three approaches. The hash-set implementation yields almost zero CES detection, since that pattern is not covered by the 4-entry Java API mapping table, and thus requires a WL-only retrieval. Both cases are documented as taxonomy coverage issues in Section~\ref{subsec:threats}.

\subsection{APM Generation Results}
\label{subsec:apm_results}

The APM generation module was tested separately using a fully directed 120-pair test matrix. As a supplementary, architecturally separate feasibility proof, APM yields 99.2\% global compilation success (119/120) and 98.3\% behaviorally congruent pairs (118/120) under test conditions constructed by the author. Neither result validates, nor is validated by, the results from the similarity pipeline test; both pipelines only share CPG representation. Table~\ref{tab:apm_results} lists all six directional pairs.

\begin{table}[ht]
\caption{APM Feasibility Demonstration Results (Author-Constructed Test Conditions). Behavioral compatibility tested against 3--5 test cases written by the author per each direction pair; these results indicate internal consistency under limited test conditions and do not represent independently validated performance boundaries.}
\label{tab:apm_results}
\centering
\begin{tabular}{|l|c|c|c|}
\hline
\textbf{Direction} & \textbf{Compiled} & \textbf{Behavioral Match} & \textbf{Notes} \\
\hline
C$\rightarrow$C++     & 20/20 & 20/20 & --- \\
C$\rightarrow$Java    & 20/20 & 20/20 & --- \\
C++$\rightarrow$C     & 20/20 & 20/20 & --- \\
C++$\rightarrow$Java  & 20/20 & 19/20 & P8 GCD failure \\
Java$\rightarrow$C    & 19/20 & 18/20 & P8 GCD failure \\
Java$\rightarrow$C++  & 20/20 & 19/20 & P8 GCD failure \\
\hline
\textbf{Total}        & \textbf{119/120} & \textbf{118/120} & \\
\hline
\textbf{Global Rate} & \textbf{99.2\%} & \textbf{98.3\%} & \\
\hline
\end{tabular}
\end{table}

Each of the three errors deals with Problem~8 (GCD, Euclidean algorithm, recursive solution). This occurs because the generator produces an incorrect ternary-operator condition corresponding to the recursive Euclidean GCD pattern \texttt{return b == 0 ? a : gcd(b, a\%b)}. The issue arises during parsing due to a conditional expression precedence bug in one specific version of the Java frontend of Joern. The sole compilation error (Java$\rightarrow$C) is caused by the very same pattern. Bidirectional closure (i.e., success for both $A\rightarrow B$ and $B\rightarrow A$) is reached in 19 out of 20 problems ($\geq$95\%). The 98.3\% and 99.2\% numbers in Table~\ref{tab:apm_results} do not constitute independently verified upper bounds on performance outside of our particular evaluation scenario since third-party fuzz testing using independent test cases is necessary prior to generalization. These results demonstrate feasibility under constrained evaluation conditions and should not be interpreted as general guarantees of semantic equivalence or production-grade transpilation correctness.

\subsection{CES Extraction Failure Analysis}
\label{subsec:error_analysis}

Table~\ref{tab:ces_failure} gives the zero-CES failure percentages. One important difference lies between the ablated C++ dataset and the curated cross-language evaluation dataset.

\begin{table}[ht]
\caption{Failure rates for the extraction of CES patterns in our datasets and languages. Note that 23.25\% refers to the uncurated ablation dataset, and the cross-language evaluation dataset has only a 0.25\% failure rate due to selecting compilable and interpretable programs with entry point identification.}
\label{tab:ces_failure}
\centering
\begin{tabular}{|l|c|c|c|}
\hline
\textbf{Dataset} & \textbf{Language} & \textbf{Zero-CES} & \textbf{Rate} \\
\hline
C++ Ablation (uncurated) & C++ & 93/400 & 23.25\% \\
\hline
Cross-language (curated) & C   & 0/130  & 0.00\% \\
Cross-language (curated) & C++ & 0/130  & 0.00\% \\
Cross-language (curated) & Java& 1/140  & 0.71\% \\
\textbf{Cross-lang. Overall} & & \textbf{1/400} & \textbf{0.25\%} \\
\hline
\end{tabular}
\end{table}

In the uncurated ablation dataset, zero-CES failures are caused by four categories: hidden helper function chain calls (34.4\%), recursive calls that go unrecognized (35.5\%), STL template fragility (17.2\%), and taxonomy limitations (12.9\%). Our heuristic-based entry point detection algorithm, which determines the 34.4\% hidden helper functions category, considers the entry point of a program with respect to our curated list of patterns: \{\texttt{main}, \texttt{Main}, \texttt{solve}, \texttt{Solve}, \texttt{run}, \texttt{execute}, \texttt{solution}, \texttt{compute}, \texttt{start}\} as well as all public static void functions with an instructor-specified signature for the course. The complete list of patterns (23 patterns) can be found in the README of our GitHub repository. When an assignment gets zero CES pattern score, our algorithm relies on the WL retrieval process. Out of 93 assignments that got zero CES scores, 90.3\% were accurately retrieved using WL fallback. Zero-CES pattern submissions do not generate any CES patterns and, thus, no pattern-level feedback is available---direct interpretability gap documented as a validity threat in Section~\ref{subsec:threats}.
